\begin{abstract}
UV seam placement is a critical yet labor-intensive step in 3D content creation, requiring artists to balance chart shape, seam concealment, and alignment with semantic and geometric features. Existing automatic methods are primarily based on per-object optimization, relying on handcrafted objectives to avoid distortion or on proxies from pretrained models to inject semantic information. However, these strategies are not always well aligned with seams used in industrial production pipelines, often resulting in layouts that deviate from artist-preferred seam patterns and practical production requirements. 
To address these limitations, we propose \textit{\methodname{}}, a generative model for UV seam generation that aligns with artist preferences and production requirements. 
Instead of depending on manually designed objectives and constraints, \methodname{} learns the distribution of per-edge seam labels from a large corpus of existing seam layouts using a flow-matching generative model.
A key challenge is that typical Transformer architectures used in flow matching models are designed for sequential representations, such as point clouds, and cannot naturally account for mesh topology. To enable mesh-native learning, we design a Mesh Transformer backbone that interleaves local graph attention over mesh edges with global self-attention across vertices, capturing both fine-grained geometric cues and long-range topological coherence. 
To further improve inference-time controllability and quality, we exploit the training-free inpainting capability of flow models for both localized seam refinement and constraint-guided seam generation.
Extensive experiments show that by learning priors from professional seam layout data, \methodname{} produces UV layouts that better align with artist-authored preferences and achieve superior perceptual quality compared with distortion-based and semantic-proxy baselines.
\end{abstract}